\item \model{Qwen3}

\paragraph*{بازتولید مقیاس فرکانس پایه با تکنیک \en{ABF}}
مدل‌های خانواده \model{Qwen3} \cite{qwen2025qwen3} (از ۰.۶ میلیارد تا ۲۳۵ میلیارد پارامتر) در فاز پیش‌آموزش اولیه از پنجره بافت ۴,۰۹۶ توکن استفاده کردند. در فاز سوم پیش‌آموزش جهت گسترش پنجره بافت به ۳۲,۷۶۸ توکن، راهبرد \textbf{تنظیم فرکانس مبنا (\en{Adjusted Base Frequency - ABF})} \cite{xiong2023effective} اعمال گردید. در این تبدیل، پایه لگاریتمی فرکانس‌ها به میزان صد برابر افزایش یافت:
\begin{equation}
    b_{\text{base}} = 10,000 \longrightarrow 1,000,000
\end{equation}
این دگرگونی فرکانس‌های زاویه‌ای $\theta_j$ را بازتنظیم می‌کند:
\begin{equation}
    \theta_j^{\text{ABF}} = 1000000^{-2(j-1)/d}
\end{equation}
طول موج متناظر با کانال $j$-ام به صورت $\lambda_j = 2\pi / \theta_j$ افزایش یافته و از وقوع چرخش‌های تکراری فاز در طول بافت‌های ۳۲K جلوگیری به عمل می‌آورد.

\paragraph*{توسعه چهاربرابری در زمان استنتاج با تلفیق \en{YaRN} و \en{DCA}}
برای رساندن طول بافت مدل به ۱۲۸K توکن در زمان ارزیابی و استنتاج:
\begin{itemize}
    \item مدل از مکانیزم \textbf{درون‌یابی فرکانسی غیریکنواخت YaRN} \cite{peng2023yarn} با فاکتور مقیاس $s = 4$ ($32\text{K} \times 4 = 128\text{K}$) بهره می‌برد. در این روش، ابعاد دارای فرکانس بالا بدون تغییر باقی مانده و تنها ابعاد فرکانس پایین تحت درون‌یابی خطی قرار می‌گیرند.
    \item به منظور مهار سرریز توجه در بافت‌های چندده‌هزار توکنی، تکنولوژی \textbf{توجه دو-تکه‌ای (\en{Dual Chunk Attention - DCA})} \cite{an2024training} تلفیق شده است که با تفکیک توجه درون‌تکه‌ای و بین‌تکه‌ای، نویز ناشی از تداخل موقعیت‌های دوردست را خنثی می‌سازد.
\end{itemize}